\section{结论}

医疗人工智能在过去四年经历了基础模型化、生成式化与多模态化三个层面的发展。Med-PaLM 2、Med-Gemini、AMIE 等通用与多模态大模型在多个静态基准上已超过医生水平，并在 OSCE、心脏专科、心理治疗、转诊等真实场景中取得了首批随机对照试验证据；LLaVA-Med、Med-Flamingo、BiomedGPT 等模型共同绘制出覆盖语言、影像、组学的多模态能力图景；Foresight 2、RETFound、MedSAM 等专科基础模型则在 EHR 时序、眼底、病理、CT、分割等子领域上证明结构先验不可替代。

这一演进并非范式间的相互替代，而是层级化的叠加复用。对比学习与双向编码器延续进入生成式阶段，基础模型骨架与多模态对齐组件又以专科底座与工具调用的形式在第三阶段被再度激活。能力扩张同样并非单一变量驱动：训练目标的创新可以在不增加底座规模的前提下解锁新能力，结构先验在小样本与长尾任务上具有数量级意义的优势，而通用化与专科化是必须同时推进的两个方向。

然而，能力曲线的持续上扬并未带动三条关键侧线的同步成熟：临床协作的非线性效应、评估方法学的拓扑滞后，以及训练--治理链路的隐性脆弱，共同构成下一阶段亟待补齐的结构性缺口。医疗人工智能的发展图景已经初具规模，但距离“可用、可信、可监管”的目标仍有相当距离。缩小这一距离的关键，不再仅取决于模型本身，而在于围绕模型构建的系统。
